\scanpage{92}

% L1
\uline{Coordonnées universelles} \inserted{$(n\geqslant1)$}

% L1–4
On prend. $C=V(h_0,\ldots h_n)$ comme hyperplan à l'$\infty$, \deleted{alors les $s_i$ \unsure{$\in E$|$(i\leq)$}} donc $E=X\smallsetminus C$ comme fibré affine de rg $n+1$, les $s_i$ sont donc des sections de $E$ formant une base affine

% L5–7
\deleted{et de \unsure{même} les $s_i$\unsure{$-h_0$|$-s_0$} \unsure{les} forment une base du fibré vectoriel des transl. de $E$ en termes de choix. On choisit comme base $s_0,s_1\ldots s_n$}

% L8–10
Le vectoriel $\unsure{\mathcal E|E}$ ext. de $\mathcal O_S$ par \deleted{\unsure{$\mathcal O_S$}} l'esp. des tr. $V$ de $E$ \deleted{\unsure{par une}} \inserted{enveloppe} vect. de $E$ s'identifie donc à $\mathcal O_S^{n+2}$, au vecteur \deleted{$(\lambda_0,\lambda_1,\ldots,\lambda_{n+1},\unsure{\lambda_{n+2}})$ correspond} \inserted{$(\lambda_0=1,\lambda_1,\ldots,\lambda_{n+1})$ correspond}

% L11
\[
\deleted{s_0+\lambda_1(\underbrace{s_1-s_0}_{e_0})+\lambda_2(\underbrace{s_2-s_0}_{e_1})+\ldots+\lambda_{n}(\underbrace{s_n-s_0}_{e_n})}
\]

% L12–13
\[
\begin{aligned}
\mathcal E=k s_0\oplus k s_1\oplus\ldots\oplus k s_{n+1}&\xrightarrow{\unsure{\pi|\Phi}}\mathcal O_S\\
\sum\lambda_i s_i&\longmapsto\sum\lambda_i
\end{aligned}
\]

% L14–15
$\unsure{\tau_i|\mathrm T_i}$ \inserted{\unsure{ps.|pso.|ps.-}}réflexions de $\mathcal E$ \inserted{$0\leq i\leq n$} satisfaisant

% L16
\[
\deleted{\tau_i\tau_j=\tau_j\tau_i\ \text{si }i\ne j}\qquad \tau_i=\mathrm{id}+a'_i\otimes a_i
\]

% L17–20, bracketed conditions
\[
\begin{aligned}
\deleted{\tau_i s_i=}\qquad\text{\textcircled{1}}\quad&
 \langle\deleted{\unsure{a}}s_i,a'_j\rangle=0,\quad\text{si }\deleted{i,j\ \text{non cons.}}\ \inserted{0\leq i<j\leq n},\ \deleted{j\ne i+1}\\
\text{\textcircled{3}}\quad&
 \deleted{\langle s_0,a'_j\rangle=0\quad\text{si }j\geq1}\ 
 \inserted{\langle a_j,a'_i\rangle=0\quad\text{si }0\leq i\unsure{<j|,j|\leq j}\leq n,\ j\ne i+1}\\
\deleted{\text{\textcircled{3}}\quad}&
\deleted{\langle s_0,a'_0\rangle,\langle a_0,a'_1\rangle,\ldots,\langle a_{n-1},a'_n\rangle\ \text{inversibles}}\\
\text{\textcircled{2}}\quad& s_{i+1}=\tau_i s_i\quad0\leq i\leq n\deleted{\unsure{+1}}
\end{aligned}
\]

% L18, left
% slanted
\margin{moins importante $\longleftrightarrow$}

% L21
\[
\underline{a'_j=\lambda_{jj}s'_j+\lambda_{j,j+1}s'_{j+1}\ldots+\lambda_{j,n+1}s'_{\unsure{j+1|n+1}}}\qquad\text{exprime (1) :}
\]

% L22–23
\[
\begin{aligned}
(\tau_i s_i=)\ s_i+\lambda_{ii}a_i&=s_{i+1}\qquad\deleted{a_i=}\lambda_{ii}a_i=s_{i+1}-s_i\\
&\text{ops }\lambda_{ii}=1,\ \text{donc}
\end{aligned}
\]

% L24–25
\[
\begin{aligned}
\boxed{a_i=s_{i+1}-s_i}\qquad\lambda_{ii}=1\quad &\text{i.e.}\quad a'_i=\unsure{s'_i|\beta'_i|s_i}+\lambda_{i,i+1}s'_{i+1}\ldots+\lambda_{i,n+1}s'_{n+1}\\
&\text{expriment (1) et (2)}
\end{aligned}
\]

% L26–27
\[
\left.\begin{aligned}
\lambda_{i,i+2}=\lambda_{i,i+3}=\ldots&=\lambda_{i,\deleted{\unsure{n+}}n+1}
 \quad(\unsure{\text{soit}}\ 1+\unsure{\mu_i|\alpha_i})\\
\lambda_{i,i+1}&=\lambda_i
\end{aligned}\right|\quad\text{exprime \uline{(3)}}
\]

% archive sheet number: 91
